% --- Archon physics formalization source begin ---
% archon:physics
% archon:covers PhyXMiniProblems/problem_phyx_mini_0801.lean
% archon:source-report reports/phyx_mini/problem_phyx_mini_0801.source.json

\chapter{Physics problem phyx\_mini\_0801}
\label{ch:PhyXMiniProblems_problem_phyx_mini_0801}

\paragraph{Problem source.}
\#\# Physical scenario

An iceboat is at rest on a frictionless horizontal surface as shown in figure. Due to the blowing wind, 4.0 s after the iceboat is released, it is moving to the right at 6.0 m/s (about 22 km/h, or 13 mi/h).

\#\# Question

What constant horizontal force \textbackslash{}( F\_\{\textbackslash{}mathrm\{W\}\} \textbackslash{}) does the wind exert on the iceboat? The combined mass of iceboat and rider is 200 kg.

\#\# Answer choices

- A: A: 150N

- B: B: 300N

- C: C: 450N

- D: D: 600N

\#\# Figure caption (auxiliary; use the image as primary evidence)

The image depicts a person sitting on a land yacht, which is a type of vehicle with a sail that is used to travel over land. The person is wearing protective gear or clothing suitable for sailing. The yacht has a triangular sail with the text "B1" on it. In the background, there are two leafless trees, indicating a wintry or early spring setting. The overall scene suggests movement, as there is a faint depiction of motion lines. The land yacht is positioned on a smooth surface, possibly a road or flat ground.

\#\# Dataset metadata

Dynamics; Physical Model Grounding Reasoning, Multi-Formula Reasoning

\paragraph{Recorded answer/context.}
B: B: 300N

\paragraph{Figure/image path.}
/root/proposal\_for\_physic/science-mango/phyx\_mini\_run/phyx\_data/test\_image/801.png

\paragraph{Formalization target.}
create a compiling Lean file with sorry bodies at `PhyXMiniProblems/problem\_phyx\_mini\_0801.lean`. The Lean declarations must preserve the physical quantities, dimensions or dimensional roles, figure labels, governing-law hypotheses, and final relation expressed by this problem.
Use Mathlib/Physlib names found through LeanExplore where available. If a physics API is missing, introduce faithful local abstractions rather than scalar placeholder aliases.

\begin{theorem}[Physics formalization target]
\label{thm:physics:phyx_mini_0801:target}
\lean{PhyXMiniProblems.ProblemPhyXMini0801.wind_force_exerted_on_iceboat}
\uses{def:physics:phyx-mini-0801:phyxminiproblems-problemphyxmini0801-forceinnewtons, def:physics:phyx-mini-0801:phyxminiproblems-problemphyxmini0801-iceboatwindsetup, def:physics:phyx-mini-0801:phyxminiproblems-problemphyxmini0801-matchesproblemstatement, def:physics:phyx-mini-0801:phyxminiproblems-problemphyxmini0801-matchessuppliedfigure, def:physics:phyx-mini-0801:phyxminiproblems-problemphyxmini0801-hasphysicalparameters, def:physics:phyx-mini-0801:phyxminiproblems-problemphyxmini0801-satisfiesconstantwinddynamics, def:physics:phyx-mini-0801:phyxminiproblems-problemphyxmini0801-recordedanswerchoice, def:physics:phyx-mini-0801:phyxminiproblems-problemphyxmini0801-isuniquematchingdisplayedchoice}
The assigned autoformalize agent should translate this physics problem into Lean declarations in the covered file, with theorem and lemma proof bodies written as `by sorry`.
\end{theorem}
\begin{proof}
This is an autoformalization task, not a proof task. Produce statements that can later be proved by the physics prover without weakening the physical model.
\end{proof}

% --- Archon declaration topology begin ---
\section{Declaration topology}
\begin{definition}[Mass Quantity]
  \label{def:physics:phyx-mini-0801:phyxminiproblems-problemphyxmini0801-massquantity}
  \lean{PhyXMiniProblems.ProblemPhyXMini0801.MassQuantity}
  The combined physical mass of the iceboat and rider
\end{definition}
\begin{proof}
  This is the declared model definition of Mass Quantity.
\end{proof}
\begin{definition}[Duration Quantity]
  \label{def:physics:phyx-mini-0801:phyxminiproblems-problemphyxmini0801-durationquantity}
  \lean{PhyXMiniProblems.ProblemPhyXMini0801.DurationQuantity}
  A physical elapsed duration
\end{definition}
\begin{proof}
  This is the declared model definition of Duration Quantity.
\end{proof}
\begin{definition}[Velocity Quantity]
  \label{def:physics:phyx-mini-0801:phyxminiproblems-problemphyxmini0801-velocityquantity}
  \lean{PhyXMiniProblems.ProblemPhyXMini0801.VelocityQuantity}
  A signed one-dimensional horizontal velocity
\end{definition}
\begin{proof}
  This is the declared model definition of Velocity Quantity.
\end{proof}
\begin{definition}[Acceleration Quantity]
  \label{def:physics:phyx-mini-0801:phyxminiproblems-problemphyxmini0801-accelerationquantity}
  \lean{PhyXMiniProblems.ProblemPhyXMini0801.AccelerationQuantity}
  A signed one-dimensional horizontal acceleration
\end{definition}
\begin{proof}
  This is the declared model definition of Acceleration Quantity.
\end{proof}
\begin{definition}[Force Quantity]
  \label{def:physics:phyx-mini-0801:phyxminiproblems-problemphyxmini0801-forcequantity}
  \lean{PhyXMiniProblems.ProblemPhyXMini0801.ForceQuantity}
  A signed one-dimensional horizontal force
\end{definition}
\begin{proof}
  This is the declared model definition of Force Quantity.
\end{proof}
\begin{definition}[mass Readout]
  \label{def:physics:phyx-mini-0801:phyxminiproblems-problemphyxmini0801-massreadout}
  \lean{PhyXMiniProblems.ProblemPhyXMini0801.massReadout}
  \uses{def:physics:phyx-mini-0801:phyxminiproblems-problemphyxmini0801-massquantity}
  Read a mass in a selected mass unit
\end{definition}
\begin{proof}
  This is the declared model definition of mass Readout.
\end{proof}
\begin{definition}[duration Readout]
  \label{def:physics:phyx-mini-0801:phyxminiproblems-problemphyxmini0801-durationreadout}
  \lean{PhyXMiniProblems.ProblemPhyXMini0801.durationReadout}
  \uses{def:physics:phyx-mini-0801:phyxminiproblems-problemphyxmini0801-durationquantity}
  Read a duration in a selected time unit
\end{definition}
\begin{proof}
  This is the declared model definition of duration Readout.
\end{proof}
\begin{definition}[velocity Readout]
  \label{def:physics:phyx-mini-0801:phyxminiproblems-problemphyxmini0801-velocityreadout}
  \lean{PhyXMiniProblems.ProblemPhyXMini0801.velocityReadout}
  \uses{def:physics:phyx-mini-0801:phyxminiproblems-problemphyxmini0801-velocityquantity}
  Read signed velocity in coherent selected length and time units
\end{definition}
\begin{proof}
  This is the declared model definition of velocity Readout.
\end{proof}
\begin{definition}[acceleration Readout]
  \label{def:physics:phyx-mini-0801:phyxminiproblems-problemphyxmini0801-accelerationreadout}
  \lean{PhyXMiniProblems.ProblemPhyXMini0801.accelerationReadout}
  \uses{def:physics:phyx-mini-0801:phyxminiproblems-problemphyxmini0801-accelerationquantity}
  Read signed acceleration in coherent selected length and time units
\end{definition}
\begin{proof}
  This is the declared model definition of acceleration Readout.
\end{proof}
\begin{definition}[force Readout]
  \label{def:physics:phyx-mini-0801:phyxminiproblems-problemphyxmini0801-forcereadout}
  \lean{PhyXMiniProblems.ProblemPhyXMini0801.forceReadout}
  \uses{def:physics:phyx-mini-0801:phyxminiproblems-problemphyxmini0801-forcequantity}
  Read signed force in coherent selected mass, length, and time units
\end{definition}
\begin{proof}
  This is the declared model definition of force Readout.
\end{proof}
\begin{definition}[mass In Kilograms]
  \label{def:physics:phyx-mini-0801:phyxminiproblems-problemphyxmini0801-massinkilograms}
  \lean{PhyXMiniProblems.ProblemPhyXMini0801.massInKilograms}
  \uses{def:physics:phyx-mini-0801:phyxminiproblems-problemphyxmini0801-massquantity, def:physics:phyx-mini-0801:phyxminiproblems-problemphyxmini0801-massreadout}
  SI kilogram readout of a physical mass
\end{definition}
\begin{proof}
  This is the declared model definition of mass In Kilograms.
\end{proof}
\begin{definition}[duration In Seconds]
  \label{def:physics:phyx-mini-0801:phyxminiproblems-problemphyxmini0801-durationinseconds}
  \lean{PhyXMiniProblems.ProblemPhyXMini0801.durationInSeconds}
  \uses{def:physics:phyx-mini-0801:phyxminiproblems-problemphyxmini0801-durationquantity, def:physics:phyx-mini-0801:phyxminiproblems-problemphyxmini0801-durationreadout}
  SI second readout of a physical duration
\end{definition}
\begin{proof}
  This is the declared model definition of duration In Seconds.
\end{proof}
\begin{definition}[velocity In Meters Per Second]
  \label{def:physics:phyx-mini-0801:phyxminiproblems-problemphyxmini0801-velocityinmeterspersecond}
  \lean{PhyXMiniProblems.ProblemPhyXMini0801.velocityInMetersPerSecond}
  \uses{def:physics:phyx-mini-0801:phyxminiproblems-problemphyxmini0801-velocityquantity, def:physics:phyx-mini-0801:phyxminiproblems-problemphyxmini0801-velocityreadout}
  SI metre-per-second readout of a signed horizontal velocity
\end{definition}
\begin{proof}
  This is the declared model definition of velocity In Meters Per Second.
\end{proof}
\begin{definition}[acceleration In Meters Per Second Squared]
  \label{def:physics:phyx-mini-0801:phyxminiproblems-problemphyxmini0801-accelerationinmeterspersecondsquared}
  \lean{PhyXMiniProblems.ProblemPhyXMini0801.accelerationInMetersPerSecondSquared}
  \uses{def:physics:phyx-mini-0801:phyxminiproblems-problemphyxmini0801-accelerationquantity, def:physics:phyx-mini-0801:phyxminiproblems-problemphyxmini0801-accelerationreadout}
  SI metre-per-second-squared readout of a signed horizontal acceleration
\end{definition}
\begin{proof}
  This is the declared model definition of acceleration In Meters Per Second Squared.
\end{proof}
\begin{definition}[force In Newtons]
  \label{def:physics:phyx-mini-0801:phyxminiproblems-problemphyxmini0801-forceinnewtons}
  \lean{PhyXMiniProblems.ProblemPhyXMini0801.forceInNewtons}
  \uses{def:physics:phyx-mini-0801:phyxminiproblems-problemphyxmini0801-forcequantity, def:physics:phyx-mini-0801:phyxminiproblems-problemphyxmini0801-forcereadout}
  SI newton readout of a signed horizontal force
\end{definition}
\begin{proof}
  This is the declared model definition of force In Newtons.
\end{proof}
\begin{definition}[Horizontal Direction]
  \label{def:physics:phyx-mini-0801:phyxminiproblems-problemphyxmini0801-horizontaldirection}
  \lean{PhyXMiniProblems.ProblemPhyXMini0801.HorizontalDirection}
  The two possible senses of the horizontal axis
\end{definition}
\begin{proof}
  This is the declared model definition of Horizontal Direction.
\end{proof}
\begin{definition}[Surface Condition]
  \label{def:physics:phyx-mini-0801:phyxminiproblems-problemphyxmini0801-surfacecondition}
  \lean{PhyXMiniProblems.ProblemPhyXMini0801.SurfaceCondition}
  Idealization of the horizontal support surface
\end{definition}
\begin{proof}
  This is the declared model definition of Surface Condition.
\end{proof}
\begin{definition}[Wind Force Profile]
  \label{def:physics:phyx-mini-0801:phyxminiproblems-problemphyxmini0801-windforceprofile}
  \lean{PhyXMiniProblems.ProblemPhyXMini0801.WindForceProfile}
  Whether the wind's horizontal force is constant during the observation
\end{definition}
\begin{proof}
  This is the declared model definition of Wind Force Profile.
\end{proof}
\begin{definition}[Figure Object]
  \label{def:physics:phyx-mini-0801:phyxminiproblems-problemphyxmini0801-figureobject}
  \lean{PhyXMiniProblems.ProblemPhyXMini0801.FigureObject}
  Physical objects and cues visible in image 801.png
\end{definition}
\begin{proof}
  This is the declared model definition of Figure Object.
\end{proof}
\begin{definition}[Figure Label]
  \label{def:physics:phyx-mini-0801:phyxminiproblems-problemphyxmini0801-figurelabel}
  \lean{PhyXMiniProblems.ProblemPhyXMini0801.FigureLabel}
  Literal text labels visible in the supplied raster
\end{definition}
\begin{proof}
  This is the declared model definition of Figure Label.
\end{proof}
\begin{definition}[Iceboat Figure]
  \label{def:physics:phyx-mini-0801:phyxminiproblems-problemphyxmini0801-iceboatfigure}
  \lean{PhyXMiniProblems.ProblemPhyXMini0801.IceboatFigure}
  \uses{def:physics:phyx-mini-0801:phyxminiproblems-problemphyxmini0801-horizontaldirection, def:physics:phyx-mini-0801:phyxminiproblems-problemphyxmini0801-figureobject, def:physics:phyx-mini-0801:phyxminiproblems-problemphyxmini0801-figurelabel}
  A literal transcription of the primary image. In particular, hasCalibratedMetricScale distinguishes qualitative raster evidence from the numerical mass, time, and velocity data stated in the prose
\end{definition}
\begin{proof}
  This is the declared model definition of Iceboat Figure.
\end{proof}
\begin{definition}[Iceboat Wind Setup]
  \label{def:physics:phyx-mini-0801:phyxminiproblems-problemphyxmini0801-iceboatwindsetup}
  \lean{PhyXMiniProblems.ProblemPhyXMini0801.IceboatWindSetup}
  \uses{def:physics:phyx-mini-0801:phyxminiproblems-problemphyxmini0801-massquantity, def:physics:phyx-mini-0801:phyxminiproblems-problemphyxmini0801-durationquantity, def:physics:phyx-mini-0801:phyxminiproblems-problemphyxmini0801-velocityquantity, def:physics:phyx-mini-0801:phyxminiproblems-problemphyxmini0801-accelerationquantity, def:physics:phyx-mini-0801:phyxminiproblems-problemphyxmini0801-forcequantity, def:physics:phyx-mini-0801:phyxminiproblems-problemphyxmini0801-horizontaldirection, def:physics:phyx-mini-0801:phyxminiproblems-problemphyxmini0801-surfacecondition, def:physics:phyx-mini-0801:phyxminiproblems-problemphyxmini0801-windforceprofile, def:physics:phyx-mini-0801:phyxminiproblems-problemphyxmini0801-iceboatfigure}
  Independent quantities in the physical setup. Neither windHorizontalForce nor constantHorizontalAcceleration is defined from an answer choice or from the requested numerical result; the governing relations are stated below
\end{definition}
\begin{proof}
  This is the declared model definition of Iceboat Wind Setup.
\end{proof}
\begin{definition}[Matches Problem Statement]
  \label{def:physics:phyx-mini-0801:phyxminiproblems-problemphyxmini0801-matchesproblemstatement}
  \lean{PhyXMiniProblems.ProblemPhyXMini0801.MatchesProblemStatement}
  \uses{def:physics:phyx-mini-0801:phyxminiproblems-problemphyxmini0801-massinkilograms, def:physics:phyx-mini-0801:phyxminiproblems-problemphyxmini0801-durationinseconds, def:physics:phyx-mini-0801:phyxminiproblems-problemphyxmini0801-velocityinmeterspersecond, def:physics:phyx-mini-0801:phyxminiproblems-problemphyxmini0801-iceboatwindsetup}
  Numerical and qualitative information stated in the exercise prose. The signed values use rightward as positive. This structure contains no force value and does not identify a correct answer choice
\end{definition}
\begin{proof}
  This is the declared model definition of Matches Problem Statement.
\end{proof}
\begin{definition}[Matches Supplied Figure]
  \label{def:physics:phyx-mini-0801:phyxminiproblems-problemphyxmini0801-matchessuppliedfigure}
  \lean{PhyXMiniProblems.ProblemPhyXMini0801.MatchesSuppliedFigure}
  \uses{def:physics:phyx-mini-0801:phyxminiproblems-problemphyxmini0801-iceboatwindsetup}
  Literal and qualitative facts read from the supplied image. The raster's arrow fixes the displayed direction, while the numerical velocity still comes from the prose rather than a pixel measurement
\end{definition}
\begin{proof}
  This is the declared model definition of Matches Supplied Figure.
\end{proof}
\begin{definition}[Has Physical Parameters]
  \label{def:physics:phyx-mini-0801:phyxminiproblems-problemphyxmini0801-hasphysicalparameters}
  \lean{PhyXMiniProblems.ProblemPhyXMini0801.HasPhysicalParameters}
  \uses{def:physics:phyx-mini-0801:phyxminiproblems-problemphyxmini0801-massinkilograms, def:physics:phyx-mini-0801:phyxminiproblems-problemphyxmini0801-durationinseconds, def:physics:phyx-mini-0801:phyxminiproblems-problemphyxmini0801-iceboatwindsetup}
  Positivity and nondegeneracy of the two scalar physical parameters
\end{definition}
\begin{proof}
  This is the declared model definition of Has Physical Parameters.
\end{proof}
\begin{definition}[Satisfies Constant Acceleration Velocity Law]
  \label{def:physics:phyx-mini-0801:phyxminiproblems-problemphyxmini0801-satisfiesconstantaccelerationvelocitylaw}
  \lean{PhyXMiniProblems.ProblemPhyXMini0801.SatisfiesConstantAccelerationVelocityLaw}
  \uses{def:physics:phyx-mini-0801:phyxminiproblems-problemphyxmini0801-durationquantity, def:physics:phyx-mini-0801:phyxminiproblems-problemphyxmini0801-velocityquantity, def:physics:phyx-mini-0801:phyxminiproblems-problemphyxmini0801-accelerationquantity, def:physics:phyx-mini-0801:phyxminiproblems-problemphyxmini0801-durationreadout, def:physics:phyx-mini-0801:phyxminiproblems-problemphyxmini0801-velocityreadout, def:physics:phyx-mini-0801:phyxminiproblems-problemphyxmini0801-accelerationreadout}
  The constant-acceleration velocity law, stated in every coherent choice of length and time units. This is the generic relation v\_f = v\_i + a Δt, not the evaluated acceleration requested by the calculation
\end{definition}
\begin{proof}
  This is the declared model definition of Satisfies Constant Acceleration Velocity Law.
\end{proof}
\begin{definition}[Satisfies Newtons Second Law]
  \label{def:physics:phyx-mini-0801:phyxminiproblems-problemphyxmini0801-satisfiesnewtonssecondlaw}
  \lean{PhyXMiniProblems.ProblemPhyXMini0801.SatisfiesNewtonsSecondLaw}
  \uses{def:physics:phyx-mini-0801:phyxminiproblems-problemphyxmini0801-massquantity, def:physics:phyx-mini-0801:phyxminiproblems-problemphyxmini0801-accelerationquantity, def:physics:phyx-mini-0801:phyxminiproblems-problemphyxmini0801-forcequantity, def:physics:phyx-mini-0801:phyxminiproblems-problemphyxmini0801-massreadout, def:physics:phyx-mini-0801:phyxminiproblems-problemphyxmini0801-accelerationreadout, def:physics:phyx-mini-0801:phyxminiproblems-problemphyxmini0801-forcereadout}
  Newton's second law for signed one-dimensional components, stated in every coherent choice of mass, length, and time units. This local predicate mirrors the dimensionally correct F = m a relation without fixing any numerical force
\end{definition}
\begin{proof}
  This is the declared model definition of Satisfies Newtons Second Law.
\end{proof}
\begin{definition}[Satisfies Constant Wind Dynamics]
  \label{def:physics:phyx-mini-0801:phyxminiproblems-problemphyxmini0801-satisfiesconstantwinddynamics}
  \lean{PhyXMiniProblems.ProblemPhyXMini0801.SatisfiesConstantWindDynamics}
  \uses{def:physics:phyx-mini-0801:phyxminiproblems-problemphyxmini0801-iceboatwindsetup, def:physics:phyx-mini-0801:phyxminiproblems-problemphyxmini0801-satisfiesconstantaccelerationvelocitylaw, def:physics:phyx-mini-0801:phyxminiproblems-problemphyxmini0801-satisfiesnewtonssecondlaw}
  The governing horizontal dynamics. Frictionlessness and the absence of any other stated horizontal force make the wind force equal to the net horizontal force. Constant-force motion then obeys the constant-acceleration kinematic law and Newton's second law. No field states the requested 300 N result
\end{definition}
\begin{proof}
  This is the declared model definition of Satisfies Constant Wind Dynamics.
\end{proof}
\begin{lemma}[constant Horizontal Acceleration si]
  \label{lem:physics:phyx-mini-0801:phyxminiproblems-problemphyxmini0801-constanthorizontalacceleration-si}
  \lean{PhyXMiniProblems.ProblemPhyXMini0801.constantHorizontalAcceleration_si}
  \uses{def:physics:phyx-mini-0801:phyxminiproblems-problemphyxmini0801-accelerationinmeterspersecondsquared, def:physics:phyx-mini-0801:phyxminiproblems-problemphyxmini0801-iceboatwindsetup, def:physics:phyx-mini-0801:phyxminiproblems-problemphyxmini0801-matchesproblemstatement, def:physics:phyx-mini-0801:phyxminiproblems-problemphyxmini0801-hasphysicalparameters, def:physics:phyx-mini-0801:phyxminiproblems-problemphyxmini0801-satisfiesconstantwinddynamics}
  The prose readouts and the constant-acceleration velocity law imply the rightward acceleration a = (6 - 0) / 4 = 3/2 m/s². This is an intermediate derived result, not a source or governing-law assumption
\end{lemma}
\begin{proof}
  The conclusion follows from the governing relations and figure readouts named in the statement of constant Horizontal Acceleration si.
\end{proof}
\begin{definition}[Answer Choice]
  \label{def:physics:phyx-mini-0801:phyxminiproblems-problemphyxmini0801-answerchoice}
  \lean{PhyXMiniProblems.ProblemPhyXMini0801.AnswerChoice}
  Labels of the four force choices printed with the problem
\end{definition}
\begin{proof}
  This is the declared model definition of Answer Choice.
\end{proof}
\begin{definition}[force In Newtons]
  \label{def:physics:phyx-mini-0801:phyxminiproblems-problemphyxmini0801-answerchoice-forceinnewtons}
  \lean{PhyXMiniProblems.ProblemPhyXMini0801.AnswerChoice.forceInNewtons}
  \uses{def:physics:phyx-mini-0801:phyxminiproblems-problemphyxmini0801-forceinnewtons, def:physics:phyx-mini-0801:phyxminiproblems-problemphyxmini0801-answerchoice}
  Force in newtons printed beside each displayed answer label
\end{definition}
\begin{proof}
  This is the declared model definition of force In Newtons.
\end{proof}
\begin{definition}[recorded Answer Choice]
  \label{def:physics:phyx-mini-0801:phyxminiproblems-problemphyxmini0801-recordedanswerchoice}
  \lean{PhyXMiniProblems.ProblemPhyXMini0801.recordedAnswerChoice}
  \uses{def:physics:phyx-mini-0801:phyxminiproblems-problemphyxmini0801-answerchoice}
  The answer label recorded in the supplied dataset metadata
\end{definition}
\begin{proof}
  This is the declared model definition of recorded Answer Choice.
\end{proof}
\begin{definition}[Matches Displayed Force]
  \label{def:physics:phyx-mini-0801:phyxminiproblems-problemphyxmini0801-matchesdisplayedforce}
  \lean{PhyXMiniProblems.ProblemPhyXMini0801.MatchesDisplayedForce}
  \uses{def:physics:phyx-mini-0801:phyxminiproblems-problemphyxmini0801-forcequantity, def:physics:phyx-mini-0801:phyxminiproblems-problemphyxmini0801-forceinnewtons, def:physics:phyx-mini-0801:phyxminiproblems-problemphyxmini0801-answerchoice}
  Exact agreement of a physical signed force with a displayed choice
\end{definition}
\begin{proof}
  This is the declared model definition of Matches Displayed Force.
\end{proof}
\begin{definition}[Is Unique Matching Displayed Choice]
  \label{def:physics:phyx-mini-0801:phyxminiproblems-problemphyxmini0801-isuniquematchingdisplayedchoice}
  \lean{PhyXMiniProblems.ProblemPhyXMini0801.IsUniqueMatchingDisplayedChoice}
  \uses{def:physics:phyx-mini-0801:phyxminiproblems-problemphyxmini0801-forcequantity, def:physics:phyx-mini-0801:phyxminiproblems-problemphyxmini0801-answerchoice, def:physics:phyx-mini-0801:phyxminiproblems-problemphyxmini0801-matchesdisplayedforce}
  A displayed choice is the unique exact match for a physical force
\end{definition}
\begin{proof}
  This is the declared model definition of Is Unique Matching Displayed Choice.
\end{proof}
% --- Archon declaration topology end ---
% --- Archon physics formalization source end ---
